\chapter*{Summary}
The \emph{quantum many-body problem} has taken centre stage in theoretical physics for the better part of the last century. Its central challenge is to predict and understand the \emph{phase diagram} of a microscopic \emph{quantum Hamiltonian} describing a macroscopic number of interacting degrees of freedom. Due to the \emph{exponential scaling} of the dimension of the many-body \emph{Hilbert space}, it is typically unfeasible to obtain analytical solutions for more than a handful of particles. Progress on this front therefore calls for \emph{efficient} and \emph{scalable} \emph{approximate} parametrisations of the wave functions of interest.

Having emerged from the field of \emph{quantum information theory}, and in particular \emph{entanglement theory}, \emph{tensor networks} have gained significant attention over the last twenty years in the study of local quantum many-body Hamiltonians in low spatial dimensions. Tensor networks parametrise the class of wave functions satisfying an \emph{area law} scaling of their \emph{entanglement entropy}, which is observed to hold for local \emph{gapped} Hamiltonians. Their central idea is that \emph{global} properties of the wave function can be efficiently compressed into \emph{local} tensors that contain the variational parameters of the wave function. One of the main merits of this local parametrisation is that the presence of \emph{symmetries} translates into local invariance properties of the constituent tensors.

\emph{Symmetry}, and in particular \emph{symmetry breaking}, has long been appreciated as an organising principle for classical phase transitions, before becoming equally central to the study of quantum many-body phase diagrams. Traditionally, symmetries act on the underlying quantum degrees of freedom by unitary or anti-unitary representations of \emph{groups}. Within the \emph{Landau paradigm} of phase transitions, a quantum phase of matter is then characterised by the \emph{unbroken} symmetry of the ground state subspace.

In recent years, however, it has been recognised that there exist phases that go beyond the traditional Landau paradigm. These phases are characterised by the existence of \emph{hidden} symmetry operators that act on the entanglement patterns of the ground state. Since tensor networks provide direct access to these entanglement degrees of freedom, they reveal these hidden symmetries and form a natural tool for their study. Remarkably, these hidden symmetry operators need not be invertible and may fall outside the scope of group theory. Rather, they are naturally described by \emph{(higher-)fusion categories}. In a sense, tensor networks provide concrete \emph{representations} of these generalised symmetry structures.

These \emph{categorical} symmetries have also been recognised to be inherited by the \emph{edge} theories of these beyond-Landau phases. The edge theories therefore live in one lower dimension. This \emph{holographic} way of describing generalised notions of symmetry has gained significant attention in recent years and forms one of the central topics of this thesis. Concretely, this modern perspective on global symmetries allows us to use established tools from \emph{topological quantum field theory} in their study. This point of view facilitates in particular the study of \emph{gauging} and \emph{dualities} of these symmetries, as well as the classification of phases protected by them, thereby fitting these phases into a \emph{generalised Landau paradigm}.

\bigskip\noindent This thesis is about various aspects of (generalised) symmetry, in particular their associated dualities and anomalies, within the context of tensor networks. It is divided into two parts, \hyperref[part:Foundations]{\emph{Foundations}} and \hyperref[part:Applications]{\emph{Applications}}, that respectively contain introductory chapters to the topics of this thesis and the novel research results in the form of five publications, as well as a set of lecture notes.

In the introductory chapter~\ref{chap:introduction}, we begin by describing in detail the challenge posed by the quantum many-body problem. We recapitulate some key facts about the traditional Landau paradigm and the role that symmetry plays in it.

Chapter~\ref{ch:TN} is devoted to the introduction of tensor network states and some of their properties. Both the one-dimensional matrix product states (MPSs) and the two-dimensional projected entangled-pair states (PEPSs) are discussed. By making use of the fundamental theorem of matrix product states --- which stipulates in particular that an MPS is invariant under a certain symmetry if and only if its local tensors are left invariant by the symmetry --- we review in detail the classification of one-dimensional gapped phases of matter. These phases are fully described by the pattern of symmetry breaking in their ground state subspace, together with a topological invariant for their residual symmetry, representing the symmetry-protected topological order.

In \cref{ch:topo}, we first provide an introduction to the necessary notions of fusion category theory. Next, we revisit the construction of Levin and Wen's string-net models. These provide a framework for lattice realisations of non-chiral (2+1)-dimensional topological orders with gapped boundary conditions. We discuss in particular the classification of those gapped boundary conditions, as well as their ground states and anyonic excitation spectrum. Crucially, the ground states of these string-net models admit exact tensor network representations that are characterised by topological symmetries that act only on the entanglement degrees of freedom. These form precisely the hidden symmetries mentioned above.

Finally, in \cref{ch:GenSym} we construct general boundary Hamiltonians for the string-net models. These Hamiltonians can be either gapped or gapless, but inherit the (generically non-invertible) symmetries from the bulk. We then sketch how these symmetries can be gauged or dualised, and illustrate this via the Kramers-Wannier duality of the transverse field Ising model. When defined on periodic chains, duality transformations act as non-invertible linear operators. They can be promoted to genuine isometries by enlarging the Hilbert space with additional degrees of freedom that implement symmetry-twisted boundary conditions, on which the duality operators can act non-trivially. Subsequently, we study the action of the two-dimensional Kramers-Wannier duality and argue that it leads to dual models with a symmetry structure described naturally by a fusion 2-category. After outlining the relevant algebraic framework, we conclude with an outlook towards a general theory of dualities in 2+1 dimensions.

The first chapter of the second part consists of lecture notes on tensor networks. This is the only application chapter that does not contain fundamentally new results. Moreover, parts of the introductory chapters are based on relevant material from these lecture notes.

The first publication, presented in \cref{app:frieze}, classifies one-dimensional bosonic phases of matter protected by quasi-one-dimensional spatial symmetries. These symmetries correspond to the so-called frieze symmetry groups. The chapter uses techniques very similar to those discussed in \cref{ch:TN} and provides, among other results, canonical forms for the phases that arise in this classification. We also give a practical algorithm to compute group cohomology and representative cocycles.

The publication presented in chapter~\ref{app:SN} establishes the result that different Morita equivalent string-net models realise the same non-chiral topological phase. Although these string-net models are microscopically very distinct, we construct an explicit constant-depth quantum circuit that maps one model to the other. This circuit acts not only on the ground state subspace, but also on the excited states when supplemented with suitable auxiliary degrees of freedom.

Chapter~\ref{app:gauging} then demonstrates that the one-dimensional duality transformations introduced in \cref{ch:GenSym} can be interpreted as gauging a particular collection of symmetries encoded in the symmetry fusion category. In particular, we demonstrate that the MPO duality intertwiners are retrieved from a gauging map, acted upon by a unitary constant-depth quantum circuit that disentangles the gauge degrees of freedom. This generalises an earlier construction due to Haegeman et al. for group symmetries. We also discuss the role of anomalies in this framework, and determine under which conditions this gauging procedure can be viewed as implementing a local version of the symmetry.

In \cref{app:codes}, we build on an earlier construction by Jos\'e Garre-Rubio. By starting from an abelian symmetry in arbitrary dimension and iteratively gauging both the symmetry and its dual symmetry, we obtain states in one dimension higher. We show that the states obtained in this way are naturally stabiliser codes, and provide numerous examples fitting into this framework.

Finally, in \cref{app:TwistedGauging}, we turn to (2+1)-dimensional dualities. We discuss a duality transformation in both one and two spatial dimensions, that amounts to the composition of a gauging transformation followed by a twisted gauging transformation. This project solves the open problem of constructing (2+1)-dimensional self-dual models, meaning models that are left invariant under a particular duality transformation. We also lift these duality transformations to unitary transformations by taking into account the appropriate boundary conditions.